\answer{ex:05:1}{$\ell\approx17\um$; $\approx100$~dias: a difusão pelo cérebro é feita pela ramificação do fio; a difusão molecular só espalha cada local de liberação por $\sim20\um$.}
\answer{ex:05:2}{$r=ab/(1+G)$, $S=1/(1+G)$ exatamente; $+1.7\,\%$ em vez de $+20\,\%$.}
\answer{ex:05:3}{$T^*=1.21\,\text{s}$ para $G=2$ e $0.19\,\text{s}$ para $G=9$: na prática $G\sim2$--$4$, redução de $3$--$5$ vezes.}
\answer{ex:05:4}{$\mathrm{CV}=1/\sqrt{2\lambda\tau_c}\approx9\cdot10^{-4}$ com taxa total $\lambda$; um único neurônio precisa de $\tau_c=12.5\,\text{s}$.}
\answer{ex:05:5}{$c\approx0.40$ nos dois casos; $\mathrm{CV}=0.050$ contra $0.158$, $\sqrt{10}$ vezes menor. As taxas dos neurônios individuais continuam proporcionais a $a_i$: só a soma é regulada.}
\answer{ex:05:6}{$N_1=\overline{A\vee B}$, $N_2=\overline{A\vee N_1}$, $N_3=\overline{B\vee N_1}$, $N_4=\overline{N_2\vee N_3}$, XOR $=N_5=\overline{N_4}$; cada comutação leva $\tau_c\ln2$, um caminho de $4$ portas leva $2.8\,\text{s}$ por XOR.}
\answer{ex:05:7}{(a)~$\tau_\gamma=6/\ln3\approx5.5\,\text{s}$. (b)~$\bar D<4\,\text{s}$ contra $D<2.2\,\text{s}$: a imprevisibilidade do atraso torna mais paciente.}
\answer{ex:05:8}{$k_2=1/\tau_d=10\,\text{s}^{-1}$, $k_1=1/\tau_d-1/\tau_\gamma=9.5\,\text{s}^{-1}$; $\tau_\gamma=1/(k_2-mk_1)$: dobra com $+2.6\,\%$ (e com $+5.3\,\%$ o horizonte fica infinito).}
